% from main.tex
%%%%%%%%%%  scan idx 283  |  folio 270  %%%%%%%%%%
localement constante, ce qui achève la preuve de (ii).

(iii) On peut supposer \(E\) strictement parfait et \(E^i=0\) pour \(i<p\) (resp. \(p-1\)). Dire que le foncteur \(M\mapsto H^p(E_x\otimes_{\mathcal O_{X,x}}^{\mathrm L}M)\) est exact à droite (resp. à gauche) sur \(\Mod(x)\) équivaut, d'après (5.2 (iii)), à dire que \(Z'^{p+1}(E_x)\) (resp. \(Z'^p(E_x)\)) (avec les notations de (ii)) est plat, donc libre de type fini (puisque de présentation finie). Comme on a \(Z'^{p+1}(E_x)\simeq Z'^{p+1}(E)_x\) (resp. \(Z'^p(E_x)\simeq Z'^p(E)_x\)), l'argument habituel (cf. EGA \(\mathrm O_I\) 5.2.7) montre qu'il existe un objet \(x\)-ponctué \(U\) de \(X\) tel que \(Z'^{p+1}(E\,|\,U)\) (resp. \(Z'^p(E\,|\,U)\)) soit libre de type fini, et l'on gagne.

\textbf{Corollaire 5.6} (EGA III 7.6.9). Soient \(X\) un schéma, \(p\in\mathbb Z\), et \(E\in\Ob D^{-}(X)\) un complexe \((p-1)\)-pseudo-cohérent. Alors :
\begin{enumerate}[label=(\roman*)]
\item Pour \(x\in X\), \(H^p(E\otimes_X^{\mathrm L}k(x))\) est de rang fini sur \(k(x)\) et la fonction
\[
 x\longmapsto\operatorname{rg}_{k(x)}H^p(E\otimes_X^{\mathrm L}k(x))
\]
est constructible et semi-continue supérieurement sur \(X\).
\item Supposons \(X\) quasi-séparé et \(E\) à cohomologie quasi-cohérente et localement bornée. Si le foncteur \(M\mapsto H^p(E\otimes^{\mathrm L}M)\) est exact sur la catégorie des \(\mathcal O_X\)-Modules quasi-cohérents, la fonction \(x\mapsto\operatorname{rg}_{k(x)}H^p(E\otimes_X^{\mathrm L}k(x))\) est localement constante et \(H^p(E)\) est localement libre de type fini. Inversement, si \(X\) est localement noethérien et réduit et si la fonction précédente est localement constante, le foncteur \(M\mapsto H^p(E\otimes^{\mathrm L}M)\) est exact sur la catégorie des \(\mathcal O_X\)-Modules quasi-cohérents.
\end{enumerate}

\emph{Preuve.} Si \(X\) est intègre de point générique \(x\), le foncteur \(M\mapsto H^p(E_x\otimes_{\mathcal O_{X,x}}^{\mathrm L}M)\) est exact sur la catégorie des \(\mathcal O_{X,x}\)-modules (puisque \(\mathcal O_{X,x}=k(x)\)), donc (5.5 (iii)), il existe un voisinage ouvert \(U\) de \(x\) tel que le foncteur \(M\mapsto H^p(E\,|\,U\otimes^{\mathrm L}M)\) soit exact sur \(\Mod(U)\), et, sur \(U\), la fonction \(y\mapsto\operatorname{rg}_{k(x)}H^p(E\otimes_X^{\mathrm L}k(y))\) est
